\documentclass[11pt,a4paper]{article}

\usepackage[T1]{fontenc}
\usepackage[utf8]{inputenc}
\usepackage[polish]{babel}
\usepackage{lmodern}
\usepackage{geometry}
\usepackage{amsmath}
\usepackage{amssymb}
\usepackage{booktabs}
\usepackage{array}
\usepackage{hyperref}
\usepackage{enumitem}
\usepackage{graphicx}
\usepackage{xcolor}

\geometry{margin=2.5cm}
\hypersetup{
    colorlinks=true,
    linkcolor=blue!50!black,
    urlcolor=blue!50!black,
    pdftitle={Crack Path Analysis w FRASTA-toolbox},
    pdfauthor={OpenAI Codex},
    pdfsubject={Opis metod i implementacji},
    pdfkeywords={FRASTA, crack path, tortuosity, orientacja, kontakt, LaTeX}
}

\title{Crack Path Analysis w FRASTA-toolbox\\
\large Opis metody, implementacji i interpretacji wynikow}
\author{}
\date{\today}

\begin{document}

\maketitle
\tableofcontents
\newpage

\section{Cel modulu}

Modul \texttt{Crack Path Analysis} sluzy do ilosciowego opisu geometrii
frontu lub sciezki propagacji pekniecia na podstawie dwoch dopasowanych
powierzchni przelomu zapisanych jako regularne mapy wysokosci.

W aktualnym stanie implementacji modul pozwala na:

\begin{itemize}[leftmargin=1.5em]
    \item wyznaczenie obszaru \emph{open/contact} na podstawie mapy roznicy,
    \item ekstrakcje sciezki pekniecia dwiema metodami,
    \item obliczenie globalnej tortuosity,
    \item obliczenie lokalnej tortuosity w oknie przesuwanym,
    \item analize lokalnej krzywizny trajektorii,
    \item analize orientacji stycznej do trajektorii,
    \item analize stabilnosci wyniku w funkcji progu separacji.
\end{itemize}

Modul ten nie zastępuje interpretacji mechanicznej. Dostarcza natomiast
powtarzalny zestaw deskryptorow geometrycznych, ktore moga byc dalej
porownywane z warunkami procesu, mikrostruktura materialu lub wynikami testow
wytrzymalosciowych.

\section{Dane wejsciowe i zalozenia}

\subsection{Model danych}

Kazda powierzchnia jest reprezentowana jako regularna dwuwymiarowa siatka
wysokosci:

\begin{itemize}[leftmargin=1.5em]
    \item \texttt{height} -- macierz wartosci wysokosci,
    \item \texttt{dx}, \texttt{dy} -- krok probkowania w osi \(X\) i \(Y\),
    \item jednostka wewnetrzna -- mikrometry.
\end{itemize}

Modul zaklada, ze:

\begin{enumerate}[leftmargin=1.5em]
    \item obie powierzchnie zostaly juz wzajemnie dopasowane,
    \item obie siatki maja zgodny rozmiar,
    \item brakujace lub niewazne punkty sa reprezentowane przez \texttt{NaN}.
\end{enumerate}

\subsection{Mapa roznicy}

Punktem wyjscia jest mapa roznicy:

\begin{equation}
    D(x, y) = H_{\mathrm{ref}}(x, y) - H_{\mathrm{adj}}(x, y),
\end{equation}

gdzie:

\begin{itemize}[leftmargin=1.5em]
    \item \(H_{\mathrm{ref}}\) -- powierzchnia referencyjna,
    \item \(H_{\mathrm{adj}}\) -- powierzchnia dopasowana.
\end{itemize}

Interpretacja znaku jest nastepujaca:

\begin{itemize}[leftmargin=1.5em]
    \item male wartosci \(D\) odpowiadaja obszarom kontaktu lub niemal
    kontaktu,
    \item duze wartosci \(D\) odpowiadaja obszarom otwarcia.
\end{itemize}

\subsection{Prog separacji}

Uzytkownik wybiera prog \(s\), na podstawie ktorego budowany jest binarny
obszar otwarcia:

\begin{equation}
    \mathrm{open}(x, y) =
    \begin{cases}
        1, & \text{gdy } D(x, y) \ge s, \\
        0, & \text{gdy } D(x, y) < s.
    \end{cases}
\end{equation}

Prog \(s\) jest zatem parametrem segmentacji. Nie jest on w aktualnej
implementacji wyznaczany automatycznie; jest wybierany przez uzytkownika i
powinien byc interpretowany jako kryterium analityczne, a nie uniwersalna
wlasciwosc materialu.

\section{Metody ekstrakcji sciezki}

\subsection{Metoda \texttt{first\_open\_pixel}}

Pierwsza metoda jest celowo prosta i stabilna. Dla kazdej linii prostopadlej
do nominalnego kierunku propagacji wybierany jest pierwszy piksel nalezacy do
obszaru \texttt{open}.

Kierunek propagacji moze byc zdefiniowany wzdluz osi \(X\), wzdluz osi \(Y\)
albo jako dowolny kat w plaszczyznie. Dla przypadkow osiowych wybierany jest
pierwszy lub ostatni piksel \texttt{open} w zaleznosci od ustawienia
\texttt{front side = min/max}. Dla przypadku katowego ta sama logika jest
realizowana po projekcji punktow na zadany kierunek propagacji i kierunek
poprzeczny.

Zalety tej metody:

\begin{itemize}[leftmargin=1.5em]
    \item prostota,
    \item duza powtarzalnosc,
    \item niewielka wrazliwosc na lokalne schodkowanie konturu.
\end{itemize}

Ograniczenia:

\begin{itemize}[leftmargin=1.5em]
    \item metoda narzuca jeden punkt na kazda kolumne lub wiersz,
    \item w praktyce wygładza geometrie frontu,
    \item moze pomijac drobniejsze lokalne odnogi lub zafalowania.
\end{itemize}

\subsection{Metoda \texttt{contour}}

Druga metoda sledzi granice pomiedzy obszarem \texttt{open} i \texttt{contact}
na podstawie konturu maski binarnej. Sama surowa granica pikselowa jest jednak
zbyt niestabilna, aby bezposrednio liczyc z niej krzywizne i deskryptory
lokalne. Dlatego w aktualnej implementacji zastosowano trzy dodatkowe kroki:

\begin{enumerate}[leftmargin=1.5em]
    \item \textbf{Uporzadkowanie wzgledem osi propagacji} -- kontur jest
    orientowany tak, aby wspolrzedna propagacji rosla monotonicznie.
    \item \textbf{Resampling} -- trajektoria jest probkowana do stalego kroku
    wzdłuż osi propagacji.
    \item \textbf{Wygładzanie} -- skladowa poprzeczna trajektorii jest
    wygładzana w oknie ruchomym przed obliczeniem krzywizny i orientacji.
\end{enumerate}

Zalety tej metody:

\begin{itemize}[leftmargin=1.5em]
    \item lepiej sledzi rzeczywista granice \texttt{open/contact},
    \item jest blizsza bezposredniej geometrii frontu,
    \item daje bogatsza odpowiedz lokalna.
\end{itemize}

Ograniczenia:

\begin{itemize}[leftmargin=1.5em]
    \item jest bardziej wrazliwa na jakosc segmentacji,
    \item silniej reaguje na schodkowanie siatki i drobne artefakty,
    \item wymaga sensownego doboru kroku resamplingu i wygładzania.
\end{itemize}

\section{Deskryptory globalne}

\subsection{Dlugosc rozwinięta i rzutowana}

Dla wyznaczonej polilinii \(\mathbf{p}_i = (x_i, y_i)\) liczona jest:

\begin{equation}
    L_{\mathrm{eff}} = \sum_{i=1}^{N-1} \left\| \mathbf{p}_{i+1} - \mathbf{p}_i \right\|,
\end{equation}

czyli efektywna dlugosc rozwinięta trajektorii.

Nastepnie liczona jest dlugosc rzutowana na nominalny kierunek propagacji:

\begin{equation}
    L_{\mathrm{proj}} = \max(u_i) - \min(u_i),
\end{equation}

gdzie \(u_i\) jest wspolrzedna punktu w osi propagacji.

\subsection{Globalna tortuosity}

Podstawowy deskryptor globalny definiowany jest jako:

\begin{equation}
    \tau = \frac{L_{\mathrm{eff}}}{L_{\mathrm{proj}}}.
\end{equation}

Interpretacja:

\begin{itemize}[leftmargin=1.5em]
    \item \(\tau \approx 1\) -- trajektoria prawie prosta,
    \item \(\tau > 1\) -- trajektoria rozwinięta lub meandrujaca,
    \item im wieksza wartosc \(\tau\), tym bardziej zlozona jest droga
    propagacji w skali calej sciezki.
\end{itemize}

\section{Deskryptory lokalne}

\subsection{Krzywizna lokalna}

Po wyznaczeniu wspolrzednej luku \(s\) liczony jest lokalny kat stycznej
\(\theta(s)\), a nastepnie przyblizona krzywizna:

\begin{equation}
    \kappa(s) = \frac{d\theta}{ds}.
\end{equation}

Interpretacja:

\begin{itemize}[leftmargin=1.5em]
    \item wartosci bliskie zeru odpowiadaja lokalnie prostym odcinkom,
    \item duze dodatnie i ujemne wartosci odpowiadaja miejscom gwaltownej
    zmiany kierunku,
    \item srednia \(|\kappa|\) jest miara lokalnej ``ostrosci'' trajektorii.
\end{itemize}

\subsection{Lokalna tortuosity}

Lokalna tortuosity \(\tau_{\mathrm{local}}(s)\) jest liczona w oknie
przesuwanym o zadanej dlugosci \(w\) wzdluz trajektorii.

Dla kazdego polozenia okna porownywane sa:

\begin{itemize}[leftmargin=1.5em]
    \item lokalna dlugosc rozwinięta,
    \item lokalna dlugosc rzutowana.
\end{itemize}

Definicja jest analogiczna do wersji globalnej:

\begin{equation}
    \tau_{\mathrm{local}}(s) =
    \frac{L_{\mathrm{eff,local}}(s)}{L_{\mathrm{proj,local}}(s)}.
\end{equation}

W praktyce:

\begin{itemize}[leftmargin=1.5em]
    \item \(\tau_{\mathrm{local}} \approx 1\) oznacza lokalnie prawie prosty
    fragment,
    \item podwyzszone wartosci wskazuja, gdzie trajektoria jest lokalnie
    bardziej rozwinięta,
    \item parametr \texttt{Local window [\textmu m]} steruje skala usredniania.
\end{itemize}

To nie jest wykres zaleznosci od progu, lecz profil po dlugosci sciezki.
Oznacza to, ze os pozioma odpowiada polozeniu wzdluz wyznaczonej trajektorii,
a nie wartosci \(s\).

\section{Analiza orientacji}

\subsection{Kat stycznej i histogram orientacji}

Po wyznaczeniu trajektorii liczony jest lokalny kat stycznej do sciezki.
Nastepnie budowany jest histogram orientacji osiowych, to znaczy takich, dla
ktorych orientacje \(0^\circ\) i \(180^\circ\) sa traktowane jako rownowazne.

Taka reprezentacja jest naturalna dla trajektorii, w ktorej interesuje nas
kierunek geometryczny odcinka, a nie jego orientacja zwrotna.

\subsection{Dominujacy kierunek}

Modul raportuje dominujaca orientacje osiowa:

\begin{equation}
    \theta_{\mathrm{dom}} \in [0^\circ, 180^\circ),
\end{equation}

oraz sile orientacji:

\begin{equation}
    R \in [0, 1],
\end{equation}

gdzie:

\begin{itemize}[leftmargin=1.5em]
    \item \(R \approx 1\) oznacza bardzo uporzadkowany rozklad orientacji,
    \item mniejsze \(R\) oznacza szerszy i bardziej rozproszony rozklad.
\end{itemize}

\subsection{Kat referencyjny}

Pole \texttt{Ref angle} nie zmienia samej geometrii sciezki ani nie wplywa na
tortuosity. Jest to wyłącznie kat odniesienia do interpretacji orientacji.

Moze on odpowiadac na przyklad:

\begin{itemize}[leftmargin=1.5em]
    \item kierunkowi budowy (\emph{build direction}),
    \item kierunkowi hatchingu,
    \item osi probki,
    \item nominalnemu kierunkowi obciazenia.
\end{itemize}

Na tej podstawie liczona jest roznica:

\begin{equation}
    \Delta \theta_{\mathrm{ref}} =
    \min\left(
        |\theta_{\mathrm{dom}} - \theta_{\mathrm{ref}}|,
        180^\circ - |\theta_{\mathrm{dom}} - \theta_{\mathrm{ref}}|
    \right).
\end{equation}

\section{Analiza stabilnosci wyniku: \(\tau(s)\)}

Oprocz pojedynczego wyniku dla wybranego progu modul oblicza sweep:

\begin{equation}
    \tau = \tau(s)
\end{equation}

w zakresie obejmujacym aktualne wartosci mapy roznicy.

Cel tego wykresu jest praktyczny: sprawdzic, czy wynik globalnej tortuosity
jest stabilny dla rozsądnego zakresu progow.

Interpretacja:

\begin{itemize}[leftmargin=1.5em]
    \item \textbf{plateau} -- wynik stabilny i wiarygodny,
    \item \textbf{pojedynczy skok, potem plateau} -- czesto akceptowalny
    przypadek przejscia od segmentacji zbyt liberalnej do poprawnej,
    \item \textbf{liczne skoki} -- silna zaleznosc od progu,
    \item \textbf{ciagly trend bez plateau} -- brak wyraznego zakresu stabilnej
    interpretacji.
\end{itemize}

W praktyce zaleca sie wybierac prog z obszaru plateau, a nie z zakresu, w
ktorym \(\tau(s)\) zmienia sie gwaltownie.

\section{Interfejs uzytkownika}

\subsection{Powiazanie z Contact Map Analysis}

Analiza crack path jest aktualnie realizowana w osobnym oknie
\texttt{CrackPathDialog}, otwieranym z poziomu okna
\texttt{ContactMapDialog}.

Takie rozwiazanie zostalo przyjete z dwoch powodow:

\begin{enumerate}[leftmargin=1.5em]
    \item okno kontaktu pozostaje czytelne i nie jest przeciazone informacja,
    \item analiza tortuosity i orientacji moze rozwinac sie niezaleznie od
    podstawowej segmentacji \texttt{open/contact}.
\end{enumerate}

Okno crack path pozostaje powiazane z tym samym zestawem danych i moze
otrzymywac aktualizacje progu z okna kontaktu.

\subsection{Glowne kontrolki}

W aktualnym interfejsie dostepne sa miedzy innymi:

\begin{itemize}[leftmargin=1.5em]
    \item \texttt{Threshold s [\textmu m]},
    \item \texttt{Method},
    \item \texttt{Propagation direction},
    \item \texttt{Angle [deg]} dla trybu \texttt{Manual angle},
    \item \texttt{Front side},
    \item \texttt{Resample [\textmu m]} dla metody \texttt{contour},
    \item \texttt{Smooth win} dla metody \texttt{contour},
    \item \texttt{Local window [\textmu m]},
    \item \texttt{Ref angle [deg]}.
\end{itemize}

\subsection{Zakladki i wykresy}

Aktualne okno zawiera:

\begin{itemize}[leftmargin=1.5em]
    \item podglad obszaru \texttt{open} z nałożoną trajektoria,
    \item wykres krzywizny,
    \item wykres lokalnej tortuosity,
    \item wykres \(\tau(s)\),
    \item histogram orientacji.
\end{itemize}

\subsection{Eksport}

Modul pozwala eksportowac:

\begin{itemize}[leftmargin=1.5em]
    \item dane sciezki i profili do pliku \texttt{NPZ},
    \item statystyki i deskryptory do pliku \texttt{JSON}.
\end{itemize}

\section{Przyklady syntetyczne}

W katalogu \texttt{examples/data/} przygotowano kilka kontrolowanych zbiorow:

\begin{itemize}[leftmargin=1.5em]
    \item \texttt{crack\_path\_straight\_demo},
    \item \texttt{crack\_path\_wavy\_demo},
    \item \texttt{crack\_path\_y\_axis\_demo},
    \item \texttt{crack\_path\_realistic\_demo}.
\end{itemize}

Kazdy zbior ma odpowiadajacy mu plik \texttt{JSON} z sugerowanym progiem,
kierunkiem propagacji, strona frontu i wartosciami referencyjnymi. Kierunek
moze byc osiowy albo opisany jawnym katem w plaszczyznie.

Szczegolnie istotny jest zbior \texttt{realistic}, ktory zawiera:

\begin{itemize}[leftmargin=1.5em]
    \item nieregularny front,
    \item wieloskalowa chropowatosc,
    \item stopniowy wzrost otwarcia za frontem,
    \item lokalne mostki kontaktu,
    \item rzadkie braki danych.
\end{itemize}

Jest on bardziej zblizony do danych rzeczywistych niz proste przypadki
kontrolne \texttt{straight} i \texttt{wavy}.

\section{Mocne strony obecnej implementacji}

\begin{itemize}[leftmargin=1.5em]
    \item modul jest deterministyczny i testowalny na danych syntetycznych,
    \item posiada dwa warianty ekstrakcji sciezki,
    \item laczy informacje globalne, lokalne i orientacyjne,
    \item pozwala oceniac stabilnosc wyniku wzgledem progu segmentacji,
    \item jest zintegrowany z istniejacym workflow kontaktu i mapy roznicy.
\end{itemize}

\section{Ograniczenia}

Obecny stan implementacji ma kilka swiadomie przyjetych ograniczen:

\begin{itemize}[leftmargin=1.5em]
    \item brak automatycznego doboru progu,
    \item brak pelnej obslugi rozgalezien i wielu konkurencyjnych frontow,
    \item brak analizy 2D pola lokalnej tortuosity poza sama trajektoria,
    \item brak wbudowanej analizy wieloskalowej PSD lub wavelet dla trajektorii,
    \item wyniki nadal zależą od definicji frontu i ustawien segmentacji.
\end{itemize}

\section{Naturalne kierunki dalszego rozwoju}

Na podstawie dotychczasowych prac najbardziej naturalne wydaja sie:

\begin{enumerate}[leftmargin=1.5em]
    \item automatyczne wyznaczanie stabilnego zakresu progow na podstawie
    wykresu \(\tau(s)\),
    \item osobna analiza zgodnosci z \emph{build direction} i
    \emph{hatch direction},
    \item analiza wieloskalowa trajektorii,
    \item wykrywanie odnog i metryki branchingu,
    \item rozszerzenie lokalnej tortuosity do map przestrzennych 2D.
\end{enumerate}

\section{Podsumowanie}

Aktualny modul \texttt{Crack Path Analysis} w FRASTA-toolbox stanowi juz
spojny zestaw narzedzi do analizy geometrii frontu pekniecia.

W obecnym stanie uzytkownik otrzymuje:

\begin{itemize}[leftmargin=1.5em]
    \item segmentacje \texttt{open/contact},
    \item dwie metody ekstrakcji trajektorii,
    \item deskryptory globalne,
    \item deskryptory lokalne,
    \item analize orientacji,
    \item test stabilnosci wzgledem progu.
\end{itemize}

Oznacza to, ze narzedzie nie jest juz tylko wizualizacja kontaktu, lecz staje
sie pelnoprawnym modulem analitycznym, ktory mozna rozwijac w kierunku badan
materialowych, analizy SLM i bardziej zaawansowanej interpretacji zmeczeniowej.

\end{document}
